\documentclass[exercice,a4]{thedocument}%
\usepackage{texmaths-tikz}%
\usepackage{pstricks}%
\usepackage{pst-bezier}%
\usepackage{pstricks-add}%
\usepackage{pst-eucl}%
\usepackage{pst-func}%
\usepackage{pst-tree}%
\usepackage{pst-circ}%
\usepackage{tikz}%
\usepackage{pgfplots}%
\usepackage{tkz-tab}%
\usepackage{tikzpagenodes}%
\usepackage{tabularray}%
\usepackage{tabularx}%
\usepackage{multirow}%
\usepackage{multicol}%
\usepackage{scratch3}%
\usepackage{enumitem}%
\usepackage[np]{numprint}%
\usepackage{pas-tableur}%
\usepackage{listings}%
\usepackage{graphicx}%
\usepackage{adjustbox}%
\usepackage{needspace}%
\usepackage{xcolor}%
\usepackage{eurosym}%
\usepackage{diagbox}%
\usepackage{colortbl}%
\usepackage{multido}%
\usepackage{icomma}%
\usepackage[french]{babel}%
\startdatakeys
    \source{26GENMATAN1}
    \BO{2026}
    \cycle{4}
    \competencesDuSocle{CA,RA,CO}
    \objectifsApprentissage{C4A1aa,C4A2ad,C4B2cb,C4D2aa,C4D3ac}
    \motsCles{fonction,image,antécédent,tableur,graphique,équation}
\enddatakeys
\definecolor{gris}{gray}{0.5}%
\providecommand{\wideparen}[1]{\overset{\frown}{#1}}%
\begin{document}%
\providecommand{\degres}{^\circ}%
%-------------------------------------------------------------

On considère les fonctions $f$ et $g$ définies par :
$f(x) = (x-1)(x+3)$ \quad et \quad $g(x) = 2x+1$.

\begin{enumerate}
\item Calculer $f(-4)$.
\item Déterminer l'antécédent de 2 par la fonction $g$.
\item On utilise un tableur pour donner les images des nombres entiers de 0 à 8 par les fonctions $f$ et $g$.

\vspace{6pt}
\begin{center}
\begin{tabularx}{\linewidth}{|c|*{10}{>{\centering \arraybackslash}X|}}\hline
& \textbf{A} & \textbf{B} & \textbf{C} & \textbf{D} & \textbf{E} & \textbf{F} & \textbf{G} & \textbf{H} & \textbf{I} & \textbf{J} \\\hline
\textbf{1} & $x$ & 0 & 1 & 2 & 3 & 4 & 5 & 6 & 7 & 8 \\\hline
\textbf{2} & $f(x)$ & $-3$ & 0 & 5 & 12 & 21 & 32 & 45 & 60 & 77 \\\hline
\textbf{3} & $g(x)$ & 1 & 3 & 5 & 7 & 9 & 11 & 13 & 15 & 17 \\\hline
\end{tabularx}
\end{center}

\begin{enumerate}[label=\alph*)]
\item Quelle formule doit-on saisir en cellule B3 puis étirer vers la droite pour compléter la ligne 3\,?\\
\textit{Aucune justification n'est demandée.}
\item Par lecture du tableau ci-dessus, donner une solution de l'équation $f(x) = g(x)$.\\
\textit{Aucune justification n'est demandée.}
\end{enumerate}

\item On représente graphiquement chacune de ces fonctions.

\medskip

\begin{minipage}[t]{0.52\textwidth}
\begin{center}
%début Tikz
%\begin{tikzpicture}
%\begin{axis}[
%axis lines=center,
%xmin=-4.5, xmax=3.5,
%ymin=-5.5, ymax=5.5,
%xtick={-4,-3,...,3},
%ytick={-5,-4,...,5},
%tick label style={font=\small},
%grid=both,
%grid style={line width=0.3pt, draw=gray!40},
%width=7.5cm, height=8.5cm,
%]
%% C1 = courbe parabolique f(x)=(x-1)(x+3)
%\addplot[thick, domain=-4.2:3.2, samples=80] {(x-1)*(x+3)};
%\node[above left, font=\small] at (axis cs:-3.8,4.5) {$\mathcal{C}_1$};
%% C2 = droite g(x)=2x+1
%\addplot[thick, domain=-4.2:3.2, samples=2] {2*x+1};
%\node[below right, font=\small] at (axis cs:2,4) {$\mathcal{C}_2$};
%\end{axis}
%\end{tikzpicture}
%Fin Tikz
%début PSTricks
\psset{unit=1cm,arrowsize=2pt 3}
\begin{pspicture*}(-4.5,-5.5)(3.5,5.5)
\psaxes[linewidth=1.25pt,labelFontSize=\scriptstyle]{->}(0,0)(-4.5,-5.5)(3.5,5.5)
\psplot[plotpoints=2000,linewidth=1.25pt,linecolor=red]{-4}{2.5}{x 1 sub x 3 add mul}
\uput[r](-3.5,3.5){\red $\mathcal{C}_1$}
\psplot[plotpoints=200,linewidth=1.25pt]{-4}{2.5}{x 2 mul 1 add}\uput[r](-2.7,-4.5){ $\mathcal{C}_2$}
\end{pspicture*}
%fin PSTricks
\end{center}
\end{minipage}
\hfill
\begin{minipage}[t]{0.44\textwidth}
\vspace{20pt}
\begin{enumerate}[label=\alph*)]
\item Associer à chacune des fonctions $f$ et $g$ sa représentation graphique.\\
\textit{Aucune justification n'est demandée.}
\item Par lecture graphique, déterminer les deux solutions de l'équation $f(x) = g(x)$.\\
\textit{Aucune justification n'est demandée.}
\end{enumerate}
\end{minipage}

\item Lola affirme que les solutions de l'équation $f(x) = g(x)$ sont les mêmes que les solutions de l'équation $x^2 - 4 = 0$. A-t-elle raison\,? Justifier.
\end{enumerate}

\newpage

%%%%%%%%%
\medskip
\showthesource
\end{document}
